\section{11.09.2026} 

\subsection{Задача 1}

\begin{exercise}
\begin{lstlisting}
int sum = 0;

for (size_t i = 1; i < N: i *= 2) {
    for (size_t j = N; j > 0; j /= 2) {
        for (size_t k = j; k < N; k += 2) {
            sum += (i + j * k);
        }
    }
}

\end{lstlisting}

С помощью пошаговой трассировки составить точную функцию временной сложности $T(N)$ и оценить порядок роста этой функции
\end{exercise}

\begin{solution}
    \[
        \log_2 n
    \]
\begin{center}
\begin{tabular}{Sl|Sl}
$i$ & $k$\\
\hline
$N$   & $\frac{N}{2} - \frac{N}{2}$ \\
\hline
$\frac{N}{2}$   & $\frac{N}{2} - \frac{N}{4}$ \\
\hline
$\frac{N}{4}$   & $\frac{N}{2} - \frac{N}{8}$ \\
\hline
$\vdots$ & \\
\hline
$1$   & $\frac{N}{2} - \frac{1}{2}$ \\
\hline
\end{tabular}
\end{center}

\[
    S = \frac{b_1}{1-q}
\]

\[
    \log_2 N \cdot \frac{N}{2} - N \underbrace{\left( \frac{1}{2} + \frac{1}{4} + \ldots \right)}_{1}
\]
\[
    T(N) = \left( \log _2 N\cdot \frac{N}{2} - N \right) \cdot \log _2 N = \log _2^2 N \cdot \frac{N}{2} - N\log _2 N
\]
\[
    \Theta(\log ^2 N \cdot N)
\]

\end{solution}

\subsection{Задача 2}

\begin{exercise}

\begin{lstlisting}
int multiply(int a, int b) {
    int result = 0;

    for (int i = 0; i < a; ++i) {
        result += b;
    }

    return result;
}
\end{lstlisting}

\begin{itemize}
    \item Подходит ли условие $P = \{2 + 2 = 4\}$ в качестве инварианта приведенного циклического алгоритма?
    \item Какое условие является более подходящим в качестве инварианта?
\end{itemize}
    
\end{exercise}

\begin{solution}
    \[
        P = \{result = i \cdot b\} \text{ --- что вычисляется на конкретном этапе }
    \]
    \[
        result = b + b + b + \ldots = a \cdot b
    \]

    \begin{enumerate}
        \item[1)] $init : i = 0,\ result = 0 \cdot b = 0$
        \item[2)] $mnt : result = (i - 1) \cdot b = i_{old} \cdot b$
        \[
            i_{new} = i_{old} + 1,\ result = (i-1) \cdot b + b = i_{new} \cdot b  
        \]
        \item[3)] $trm i = a,\ result = a \cdot b$
    \end{enumerate}
\end{solution}

\subsection{Задача 3}

\begin{exercise}
\begin{lstlisting}
int i = 2;

while (i > 0) {
    i = i + 1;
    std::cout << i;
}

i = i - 1;
\end{lstlisting}

\begin{itemize}
    \item Подходит ли условие $P = \{i \geq 0\}$ в качестве инварианта приведенного циклического алгоритма?
    \item Как проверить выполнимость $P$ при выходе из цикла $[TRM]$?
\end{itemize}
\end{exercise}

\subsection{Логика Хоара}

\subsubsection{Тройка Хоара}

\[
    \underbrace{\{P\}}_{\text{предусловие}} S \underbrace{\{Q\}}_{\text{постусловие}}
\]

если $P$ --- истина и $S$ не завершается, то $Q$ --- истина

\subsection{Задача 4}

\begin{exercise}[Задача линейного поиска]

\textbf{Вход:}

Последовательность ключей $A = \{a_1, \ldots, a_n\}$ и значение $v$

\textbf{Выход:}

Индекс $1 \leq i \leq n$, для которого $A[i] = v$ или сообщение \textbf{NOT FOUND}, если такого индекса не существует. 

\begin{itemize}
    \item Разработайте циклический алгоритм линейного поиска
    \item Докажите корректность разработанного алгоритма с помощью инварианта
    \item Оцените сложность $T(n)$ разработанного алгоритма
\end{itemize}
\end{exercise}

\begin{solution}
\begin{lstlisting}
for i = 1 ... n:
    if (a[i] == v)
        return i 
return NOT FOUND
\end{lstlisting}

\begin{theorem}[Утверждение]
\[
    \boxed{v \in A \Leftrightarrow \exists i : A[i] = v} \text{ --- цель}
\]
\end{theorem}

\[
    P_{inv} = \{A[i] = = v \Leftrightarrow A[0] \neq v,\ A[1] \neq v, \ldots, A[i-1] \neq v\} 
\]
\end{solution}

\subsection{Задача 5}

\begin{exercise}[Задача бинарного поиска]

\textbf{Вход:}

Последовательность ключей $A = \{a_1, \ldots, a_n\}$ и значение $v$

\textbf{Выход:}

Индекс $1 \leq i \leq n$, для которого $A[i] = v$ или сообщение \textbf{NOT FOUND}, если такого индекса не существует. 

\begin{itemize}
    \item Разработайте циклический алгоритм бинарного поиска
    \item Докажите корректность разработанного алгоритма с помощью инварианта
    \item Оцените сложность $T(n)$ разработанного алгоритма
\end{itemize}
    
\end{exercise}

\begin{solution}
\begin{lstlisting}
left, right = 0, n - 1
while left <= right: 
    ind = floor((right + left) / 2)
    if A[ind] == v: 
        return ind
    elif A[ind] < v: 
        left = ind + 1
    else:
        right = ind 
return NOT FOUND  
\end{lstlisting}

\begin{theorem}[Утверждение]
    $$v \in A \Leftrightarrow \exists\ left \leq i \leq right : A[i] == v$$
\end{theorem}

\[
    T(n) = \{5 \log ^ 2 n + 1\}
\]
\end{solution}
